%% ============================================================
%%  PREAMBLE — Production-Grade Agentic AI: LangGraph + Python
%%  Extracted from main.tex.  Input ONCE from Tex/main.tex.
%% ============================================================

\documentclass[
  11pt,
  twoside,
  openright
]{book}

\usepackage{polyglossia}
\setdefaultlanguage{english}

\usepackage{fontspec}
\setmainfont{TeX Gyre Termes}[
  BoldFont   = TeX Gyre Termes Bold,
  ItalicFont = TeX Gyre Termes Italic
]
\setsansfont{TeX Gyre Heros}[
  BoldFont   = TeX Gyre Heros Bold,
  ItalicFont = TeX Gyre Heros Italic
]
\setmonofont{TeX Gyre Cursor}[
  Scale    = 0.88,
  BoldFont = TeX Gyre Cursor Bold
]

\usepackage{geometry}
\geometry{
  papersize  = {7in, 10in},
  top        = 1in,
  bottom     = 1.1in,
  inner      = 1in,
  outer      = 0.875in,
  headheight = 23.07pt,
  twoside
}

%% Fix fancyhdr \headheight warning (must be after geometry)
\setlength{\headheight}{23.07004pt}
\addtolength{\topmargin}{-0.00003pt}

\usepackage{xcolor}
\definecolor{BookBlue}       {HTML}{1A4A8A}
\definecolor{BookDark}       {HTML}{1C1C1E}
\definecolor{BookMid}        {HTML}{4A4A4A}
\definecolor{LearnBg}        {HTML}{EEF7EE}
\definecolor{LearnBorder}    {HTML}{27AE60}
\definecolor{TipBg}          {HTML}{EAF4FB}
\definecolor{TipBorder}      {HTML}{2980B9}
\definecolor{WarningBg}      {HTML}{FEF9E7}
\definecolor{WarningBorder}  {HTML}{F39C12}
\definecolor{NoteBg}         {HTML}{F4F4F4}
\definecolor{NoteBorder}     {HTML}{888888}
\definecolor{TakeawayBg}     {HTML}{F0F0FF}
\definecolor{TakeawayBorder} {HTML}{5B4FBE}
\definecolor{CodeBg}         {HTML}{F7F7F7}
\definecolor{CodeComment}    {HTML}{6A737D}
\definecolor{CodeKeyword}    {HTML}{D73A49}
\definecolor{CodeString}     {HTML}{032F62}
\definecolor{CodeConstant}   {HTML}{005CC5}
\definecolor{CodeAnnotation} {HTML}{6A737D}

\usepackage{microtype}
\usepackage{setspace}
\setstretch{1.18}

\usepackage{csquotes}
\usepackage{parskip}
\setlength{\parskip}{6pt}
\setlength{\parindent}{0pt}

\widowpenalty=10000
\clubpenalty=10000
\displaywidowpenalty=10000
\predisplaypenalty=10000
\raggedbottom
\emergencystretch=3em
\hbadness=3000
\vbadness=3000

\usepackage{needspace}
\usepackage{etoolbox}
\preto\section{\needspace{5\baselineskip}}
\preto\subsection{\needspace{4\baselineskip}}
\preto\subsubsection{\needspace{3\baselineskip}}

\usepackage{titlesec}
\renewcommand{\bottomtitlespace}{0.15\textheight}

\titleformat{\part}[display]
  {\centering\sffamily\Huge\bfseries\color{BookBlue}}
  {\sffamily\large\mdseries\partname\ \thepart}{20pt}{}

\titleformat{\chapter}[display]
  {\sffamily\huge\bfseries\color{BookBlue}}
  {\sffamily\normalsize\mdseries\chaptertitlename\ \thechapter}{12pt}{}
\titlespacing*{\chapter}{0pt}{0pt}{24pt}

\titleformat{\section}
  {\sffamily\Large\bfseries\color{BookBlue}}
  {\thesection}{1em}{}
\titlespacing*{\section}{0pt}{18pt}{6pt}

\titleformat{\subsection}
  {\sffamily\large\bfseries\color{BookDark}}
  {\thesubsection}{1em}{}
\titlespacing*{\subsection}{0pt}{14pt}{4pt}

\titleformat{\subsubsection}
  {\sffamily\normalsize\bfseries\color{BookDark}}
  {\thesubsubsection}{1em}{}
\titlespacing*{\subsubsection}{0pt}{10pt}{2pt}

\usepackage{fancyhdr}
\pagestyle{fancy}
\fancyhf{}
\fancyhead[LE]{\small\sffamily\color{BookMid}\leftmark}
\fancyhead[RO]{\small\sffamily\color{BookMid}\rightmark}
\fancyfoot[LE,RO]{\small\sffamily\color{BookMid}\thepage}
\renewcommand{\headrulewidth}{0.3pt}
\renewcommand{\footrulewidth}{0pt}

\fancypagestyle{plain}{
  \fancyhf{}
  \fancyfoot[LE,RO]{\small\sffamily\color{BookMid}\thepage}
  \renewcommand{\headrulewidth}{0pt}
}

\makeatletter
\renewcommand{\cleardoublepage}{%
  \clearpage
  \if@twoside
    \ifodd\c@page\else
      \thispagestyle{empty}\hbox{}\newpage
      \if@twocolumn\hbox{}\newpage\fi
    \fi
  \fi
}
\makeatother

\usepackage{rotating}
\usepackage{booktabs}
\usepackage{tabularx}
\usepackage{array}
\usepackage{longtable}
\renewcommand{\arraystretch}{1.35}

\newcolumntype{P}[1]{>{\RaggedRight\arraybackslash}p{#1}}

\usepackage{graphicx}
\usepackage{float}
\usepackage{caption}
\captionsetup{
  font      = {small,sf},
  labelfont = {bf,color=BookBlue},
  skip      = 6pt
}

\usepackage{listings}
\lstset{
  backgroundcolor   = \color{CodeBg},
  basicstyle        = \ttfamily\fontsize{9}{11}\selectfont,
  keywordstyle      = \color{CodeKeyword}\bfseries,
  stringstyle       = \color{CodeString}\itshape,
  commentstyle      = \color{CodeComment}\itshape,
  numbers           = left,
  numberstyle       = \tiny\color{BookMid},
  numbersep         = 6pt,
  breaklines        = true,
  frame             = single,
  framerule         = 0.4pt,
  rulecolor         = \color{NoteBorder},
  showstringspaces  = false,
  tabsize           = 2,
  captionpos        = b,
  breakatwhitespace = false,
  columns           = fullflexible,
  linewidth         = \linewidth,
  postbreak         = \mbox{\textcolor{gray}{$\hookrightarrow$}\space},
  xleftmargin       = 1.5em,
  framexleftmargin  = 1.5em,
  aboveskip         = 6pt,
  belowskip         = 4pt,
}

\lstdefinelanguage{Python}{
  morekeywords={
    False,None,True,and,as,assert,async,await,break,class,continue,def,
    del,elif,else,except,finally,for,from,global,if,import,in,is,lambda,
    nonlocal,not,or,pass,raise,return,try,while,with,yield
  },
  morekeywords=[2]{self,cls},
  sensitive=true,
  morecomment=[l]{\#},
  morestring=[b]",
  morestring=[b]',
  morestring=[s]{"""}{"""}, 
  morestring=[s]{'''}{'''},
  keywordstyle=\color{CodeKeyword}\bfseries,
  keywordstyle=[2]\color{CodeConstant},
  stringstyle=\color{CodeString},
  commentstyle=\color{CodeComment}\itshape,
  breaklines=true,
  breakatwhitespace=false,
  columns=fullflexible,
  keepspaces=true,
  showstringspaces=false,
  postbreak=\mbox{\textcolor{gray}{$\hookrightarrow$}\space},
}

\lstdefinelanguage{yaml}{
  morecomment=[l]{\#},
  morestring=[b]",
  morestring=[b]',
  literate=
    *{true}{{\color{CodeKeyword}true}}{4}
    {false}{{\color{CodeKeyword}false}}{5}
    {null}{{\color{CodeKeyword}null}}{4},
  sensitive=true,
  keywordstyle=\color{CodeKeyword}\bfseries,
  stringstyle=\color{CodeString},
  commentstyle=\color{CodeComment}\itshape,
  breaklines=true,
  columns=fullflexible,
  showstringspaces=false
}

\lstdefinelanguage{json}{
  morestring=[b]",
  morecomment=[l]//,
  morecomment=[s]{/*}{*/},
  sensitive=true,
  stringstyle=\color{CodeString},
  commentstyle=\color{CodeComment}\itshape,
}

\DeclareRobustCommand{\code}[1]{\texttt{#1}}
\providecommand{\term}[1]{\textit{#1}}
\providecommand{\chapterquote}[2]{%
  \vspace{1em}
  \begin{quote}\itshape#1\end{quote}
  \begin{flushright}--- #2\end{flushright}
  \vspace{1em}
}

\newcommand{\githubref}[1]{%
  \vspace{2pt}%
  {\small\sffamily\color{BookMid}Full source: \texttt{#1}}%
  \vspace{4pt}%
}

\usepackage{tcolorbox}
\tcbuselibrary{skins,breakable}

\newtcolorbox{learningbox}{
  colback=LearnBg, colframe=LearnBorder,
  before skip=12pt, fonttitle=\sffamily\bfseries\small,
  title={Learning Objective}, breakable, lines before break=3,
  left=8pt, right=8pt, top=6pt, bottom=6pt,
  boxrule=0.6pt, width=\linewidth
}
\newtcolorbox{tipbox}{
  colback=TipBg, colframe=TipBorder,
  before skip=12pt, fonttitle=\sffamily\bfseries\small,
  title={Tip}, breakable, lines before break=3,
  left=8pt, right=8pt, top=6pt, bottom=6pt,
  boxrule=0.6pt, width=\linewidth
}
\newtcolorbox{warningbox}{
  colback=WarningBg, colframe=WarningBorder,
  before skip=12pt, fonttitle=\sffamily\bfseries\small,
  title={Warning}, breakable, lines before break=3,
  left=8pt, right=8pt, top=6pt, bottom=6pt,
  boxrule=0.6pt, width=\linewidth
}
\newtcolorbox{notebox}{
  colback=NoteBg, colframe=NoteBorder,
  before skip=12pt, fonttitle=\sffamily\bfseries\small,
  title={Note}, breakable, lines before break=3,
  left=8pt, right=8pt, top=6pt, bottom=6pt,
  boxrule=0.6pt, width=\linewidth
}
\newtcolorbox{takeawaybox}[1][Key Takeaway]{
  colback=TakeawayBg, colframe=TakeawayBorder,
  before skip=12pt, fonttitle=\sffamily\bfseries\small,
  title={#1}, breakable, lines before break=3,
  left=8pt, right=8pt, top=6pt, bottom=6pt,
  boxrule=0.6pt, width=\linewidth
}

%% ── Source-Layout tree environment ─────────────────────────────────────
%% Usage:
%%   \begin{sourcelayout}
%%     \dirtree{%
%%       .1 src/.
%%       .2 agent/.
%%       .3 nodes/.
%%       .4 perceive\_node.py \DTcomment{§1.2}.
%%     }
%%   \end{sourcelayout}
%%
%% Requires: \usepackage{dirtree}  (loaded below)
%% Colours reuse NoteBg / NoteBorder (neutral grey — reference material,
%% not instructional content).
%% ────────────────────────────────────────────────────────────────────────
\usepackage{dirtree}

\newtcolorbox{sourcelayout}{%
  colback      = NoteBg,
  colframe     = NoteBorder,
  colbacktitle = NoteBorder,
  coltitle     = white,
  fonttitle    = \sffamily\bfseries\small,
  title        = {Source Layout},
  before skip  = 12pt,
  after skip   = 10pt,
  breakable,
  lines before break = 3,
  left   = 6pt,
  right  = 6pt,
  top    = 6pt,
  bottom = 6pt,
  boxrule = 0.6pt,
  width   = \linewidth,
}

%% dirtree global style: monospaced font, consistent with code listings
\renewcommand*\DTstylecomment{\rmfamily\color{CodeComment}\itshape}
\renewcommand*\DTstyle{\ttfamily\color{BookDark}}

\usepackage{amsmath}
\usepackage{amssymb}

\PassOptionsToPackage{hyphens}{url}
\usepackage{hyperref}
\usepackage{imakeidx}
\makeindex[columns=2, title=Index, intoc=true]
\urlstyle{tt}
\makeatletter
\g@addto@macro\UrlBreaks{\do\.\do\_\do\-}
\makeatother

\hypersetup{
  colorlinks  = true,
  linkcolor   = BookBlue,
  urlcolor    = BookBlue,
  citecolor   = BookBlue,
  pdftitle    = {Production-Grade Agentic AI: LangGraph + Python Frameworks},
  pdfauthor   = {Haythem Aber},
  pdfsubject  = {LangGraph, Python, Agentic AI, Production Systems},
  pdfkeywords = {LangGraph, Python, AI agents, production, observability, security}
}

\usepackage{tikz}
\usetikzlibrary{arrows.meta, positioning, shapes.geometric, shapes.misc, calc}
\usepackage{pdflscape}
\usepackage{ragged2e}

\newcommand{\component}[1]{\textsf{\textbf{#1}}}
\newcommand{\seechapter}[1]{\textit{(See Chapter~\ref{#1})}}
\newcommand{\seepart}[1]{\textit{(See Part~\ref{#1})}}

\newcommand{\partclose}[1]{%
  \cleardoublepage
  \vspace*{\fill}
  \begin{center}\sffamily\Large\color{BookBlue}#1\end{center}
  \vspace*{\fill}
  \cleardoublepage
}

%% Shared book-wide infrastructure layered on top of the author's preamble.
%% Keep this file loaded once by Tex/main.tex; chapter fragments must not load it.

\usepackage{calc} % Pandoc-generated longtable column dimensions.
\usepackage[
  backend=biber,
  style=verbose-note,
  autocite=footnote,
  sorting=nyt,
  maxbibnames=99
]{biblatex}
\addbibresource{Tex/bibliography/references.bib}

% Paper labels remain explicit and are not interchangeable.
\newcommand{\CoALACite}[1]{\footnote{\textsc{[CoALA]} See Sumers et al., \emph{Cognitive Architectures for Language Agents}, #1.}}
\newcommand{\FanLanCite}[1]{\footnote{\textsc{[Fan \& Lan]} See Fan and Lan, \emph{From Cognitive Architectures to Language Agents}, #1.}}

% The source paper is the theoretical basis in Parts I and II; examples and
% engineering decisions are explicitly identified in the prose.
\newcommand{\BookExampleLabel}{\textsc{[Book example]}}
\newcommand{\EngineeringLabel}{\textsc{[Engineering decision]}}

\usepackage[toc,nonumberlist]{glossaries}
\makeglossaries
\renewcommand*{\glossarypreamble}{\begin{multicols}{2}\raggedcolumns\small}
\renewcommand*{\glossarypostamble}{\end{multicols}}
\newglossaryentry{languageagent}{name={language agent},description={An agentic system that uses a language model as a core computational component and interacts with an environment}}
\newglossaryentry{cognitivelanguageagent}{name={cognitive language agent},description={A language agent that also uses internal processes to manage its information and behavior}}
\newglossaryentry{workingmemory}{name={working memory},description={Active information maintained for the current decision cycle}}
\newglossaryentry{longtermmemory}{name={long-term memory},description={Memory modules that retain information beyond the current decision cycle}}
\newglossaryentry{episodicmemory}{name={episodic memory},description={Memory of experience from earlier decision cycles}}
\newglossaryentry{semanticmemory}{name={semantic memory},description={Memory of knowledge about the world and the agent itself}}
\newglossaryentry{proceduralmemory}{name={procedural memory},description={Knowledge about how to perform actions and decision procedures, represented implicitly in model weights or explicitly in agent code}}
\newglossaryentry{grounding}{name={grounding},description={Procedures that execute external actions and process environmental feedback}}
\newglossaryentry{decisionprocedure}{name={decision procedure},description={The top-level procedure that coordinates the agent's repeated decision cycles}}


% Reusable review and recap blocks used throughout the book.
\newtcolorbox{chapterrecap}{
  colback=TakeawayBg, colframe=TakeawayBorder,
  before skip=10.5pt, fonttitle=\sffamily\bfseries\small,
  title={Chapter recap}, breakable, lines before break=3,
  left=8pt, right=8pt, top=5.25pt, bottom=5.25pt,
  boxrule=0.6pt, width=\linewidth
}
\newtcolorbox{reviewquestions}[1]{
  colback=NoteBg, colframe=NoteBorder,
  before skip=10.5pt, fonttitle=\sffamily\bfseries\small,
  title={#1}, breakable, lines before break=3,
  left=8pt, right=8pt, top=5.25pt, bottom=5.25pt,
  boxrule=0.6pt, width=\linewidth
}
\providecommand{\tightlist}{}

% Stable source references for implementation listings. The QR payload, the
% clickable source ID and "Open exact source" all consume the same URL argument.
% The visible source path/member and linked test path remain useful in print.
\newcommand{\BookSourceReferenceText}[9]{%
  {\sffamily\fontsize{8}{9}\selectfont\raggedright
    \setlength{\parskip}{0pt}%
    \href{#8}{\textbf{Source #1}}\hspace{0.55em}\textbar\hspace{0.55em}\texttt{#2}\par
    \textbf{Path:} \href{#8}{\path{#3}}\par
    \textbf{Member:} #4\hspace{0.5em}\textbf{Source lines:} #5\par
    \textbf{Test file:} \href{#9}{\path{#6}}\par
    \textbf{Tests:}\par \href{#9}{#7}\par
    \href{#8}{\textbf{Open exact source}}%
  }%
}

\newcommand{\BookSourceReferenceNoQR}[9]{%
  \par\smallskip\noindent
  \begin{minipage}{\linewidth}%
    \BookSourceReferenceText{#1}{#2}{#3}{#4}{#5}{#6}{#7}{#8}{#9}%
  \end{minipage}%
  \par\smallskip
}

\newcommand{\BookSourceReferenceQRContent}[9]{%
  \par\smallskip\noindent
  \begin{minipage}[c]{0.785\linewidth}%
    \BookSourceReferenceText{#1}{#2}{#3}{#4}{#5}{#6}{#7}{#8}{#9}%
  \end{minipage}%
  \hfill
  \begin{minipage}[c]{0.18\linewidth}%
    \centering
    % The source QR is a vector PDF asset generated with qrcode.sty.
    % Embedding it as a graphic lets hyperref create a real clickable area
    % around the full printed symbol (including its four-module quiet zone).
    \href{#8}{\mbox{{\setlength{\fboxsep}{0pt}\colorbox{white}{\includegraphics[height=\BookSourceQRHeight]{Tex/diagrams/source-qr/#1.pdf}}}}}\par
    {\scriptsize\sffamily Scan to inspect}%
  \end{minipage}%
  \par\smallskip
}

% Standard source references use a 19.5 mm QR including the quiet zone.
% Very dense commit-pinned URLs may opt into the 21 mm variant.
\newcommand{\BookSourceReferenceQR}[9]{%
  {\def\BookSourceQRHeight{19.5mm}\BookSourceReferenceQRContent{#1}{#2}{#3}{#4}{#5}{#6}{#7}{#8}{#9}}%
}
\newcommand{\BookSourceReferenceQRLong}[9]{%
  {\def\BookSourceQRHeight{21mm}\BookSourceReferenceQRContent{#1}{#2}{#3}{#4}{#5}{#6}{#7}{#8}{#9}}%
}


% Book-level metadata overrides the sample metadata retained in the supplied preamble.
\hypersetup{
  pdftitle={Reliable AI Agents: Architecture and Implementation},
  pdfauthor={Haythem Aber},
  pdfsubject={CoALA, cognitive architectures, agentic systems, Java production engineering},
  pdfkeywords={CoALA, Fan and Lan, cognitive architectures, agentic AI, Java, supply-chain disruption management, production systems}
}

\hypersetup{pdftitle={Chapter 5 Outline Review: External Actions and Grounding},pdfsubject={Detailed outline awaiting author approval},pdfkeywords={CoALA, grounding, external actions, agent environments}}
\begin{document}
\thispagestyle{empty}
\begin{center}
{\sffamily\Large\color{BookBlue} FROM COGNITIVE ARCHITECTURES TO PRODUCTION AGENTIC SYSTEMS\par}
\vspace{1.2cm}
{\sffamily\huge\bfseries\color{BookBlue} Chapter 5 Outline\par}
\vspace{0.45cm}
{\sffamily\LARGE External Actions and Grounding\par}
\vspace{0.8cm}
{\large Part I -- CoALA: Understanding Language Agents as Cognitive Systems\par}
\vfill
\begin{learningbox}
\textbf{Review status:} Awaiting author approval. This document is a detailed chapter plan, not chapter prose.

\textbf{Primary theoretical source:} Sumers et al., \emph{Cognitive Architectures for Language Agents}, especially Section 4.2. Section 4.6 is used only to place grounding inside the larger decision cycle; Section 6 is consulted briefly for the paper's discussion of action-space safety.
\end{learningbox}
\vfill
{\sffamily\large Outline review version 1.0\par}
{\small 14 November 2026\par}
\end{center}
\clearpage

\section*{1. Chapter purpose and role in the book}
Chapters 1 and 2 established CoALA's overall architecture. Chapter 3 explained what the memory categories represent, and Chapter 4 explained the internal actions that operate on those memories. The next question is what happens when an agent acts beyond its own internal information: how does it interact with a world that can respond independently of the model?

The chapter focuses on \emph{grounding}. A grounding procedure executes an external action and processes environmental feedback into working memory as text. This gives the agent a structured way to interact with physical surroundings, people or other agents, and digital environments. The teaching challenge is to help readers see that an external action is not merely a sentence generated by an LLM. It is an interaction with an environment through a procedure that can affect that environment and may produce an observation for the agent.

The chapter must not become a general survey of APIs, browser automation frameworks, robotics stacks, tool schemas, or production integration patterns. Its job is to explain CoALA's conceptual category and the three environment types described in the paper. Concrete service interactions are book-created illustrations, not theoretical claims.

\begin{takeawaybox}[Chapter thesis]
External grounding connects an agent to environments beyond its internal memories. Grounding procedures execute external actions and make environmental feedback available to the agent's working memory, allowing the agent's internal process to continue with information about what happened outside.
\end{takeawaybox}

\section*{2. Learning outcomes}
By the end of the chapter, a beginner should be able to:
\begin{enumerate}
  \item Explain why an external action is different from retrieval, reasoning, or learning.
  \item Describe grounding as a procedure that executes an external action and processes environmental feedback into working memory as text.
  \item Distinguish the environment, the grounding procedure, the action, and the observation received from the environment.
  \item Recognize the three kinds of external environments discussed by CoALA: physical environments, dialogue with humans or other agents, and digital environments.
  \item Explain how perceptual inputs can be represented as textual observations and how language-based commands can affect an external environment.
  \item Classify examples involving robots, human dialogue, websites, APIs, games, and code execution without assuming that all grounding is the same kind of tool call.
  \item State what CoALA's grounding description does not prescribe, including a universal tool protocol, a deployment architecture, or a complete safety policy.
\end{enumerate}

\section*{3. Proposed chapter structure}
\setcounter{chapter}{5}

\section{From internal work to interaction with an external world}
\textbf{Teaching entry:} Return to Product X. The agent has an active incident description and may have retrieved earlier incidents or reasoned about possible causes. But it still does not know the current state of the payment service. It must interact with something outside its own memory: for example, a monitoring endpoint, an operational dashboard, or a person who can clarify a recent change.

\textbf{Coverage:}
\begin{itemize}
  \item Recall the distinction established in Chapter 4: retrieval, reasoning, and learning interact with internal memories; grounding interacts with an external environment.
  \item Introduce the basic question: what changes or information can only be obtained by interacting with the world outside the agent?
  \item Explain that CoALA calls these external actions and groups them under grounding.
  \item Introduce an environment as the external setting with which the agent interacts; keep the definition intuitive before discussing its categories.
\end{itemize}

\textbf{Boundary:} Do not repeat the entire action-space map from Chapter 2 or the internal-action taxonomy from Chapter 4. Use those chapters as short reminders and move quickly to external interaction. \CoALACite{\S 4, Figure 5}

\section{Grounding: executing an action and processing feedback}
\textbf{Core question:} What makes an external action part of an agent's ongoing interaction rather than an isolated output?

\textbf{Coverage:}
\begin{itemize}
  \item Define grounding procedures as executing external actions and processing environmental feedback into working memory as text.
  \item Explain the two sides of the interaction: an action is sent through a grounding procedure; the environment may produce feedback or an observation that is represented for the agent.
  \item Explain why the paper describes the interaction as being simplified into a ``text game'': observations and actions can be represented textually for the language agent.
  \item Distinguish the environment from the grounding procedure. The environment is where the interaction takes place; the grounding procedure is how the agent carries out the action and processes feedback.
  \item Distinguish an action from an observation. An action is something the agent does to the environment; an observation is information obtained from the environment.
  \item Explain that observations are placed into working memory, connecting grounding back to the memory architecture without re-teaching Chapter 3.
\end{itemize}

\textbf{Planned Mermaid figure 5.1:} \emph{The grounding interaction}. Show working information leading to a grounding procedure, an external action entering an environment, and environmental feedback being represented as text and made available in working memory. Label it as a conceptual interaction, not a prescribed software pipeline.

\begin{notebox}
The chapter must preserve the paper's abstraction. It should not claim that every external action returns a successful result, that every environment responds immediately, or that every observation is naturally textual before processing. The paper describes feedback being processed into working memory as text; the implementation of that translation can vary.
\end{notebox}

\section{Physical environments: acting and perceiving in the physical world}
\textbf{Entry intuition:} A robot cannot interact with a physical room in the same way that a program calls a stateless calculator API. It must receive information from sensors and produce actions that affect the physical environment.

\textbf{Coverage:}
\begin{itemize}
  \item Explain physical embodiment as one kind of environment for an agent.
  \item Describe the perceptual side: visual, audio, and tactile inputs may be processed into textual observations, for example through pre-trained captioning models.
  \item Describe the action side: robotic planners can take language-based commands and affect the physical environment.
  \item Explain the role of vision-language models as one way to provide perceptual context to language models, without turning the chapter into a multimodal-model tutorial.
  \item Keep sensing, representation, and acting conceptually distinct. A camera image is not itself the same thing as the textual observation made available to the language model.
\end{itemize}

\textbf{Boundary:} Do not teach robotics kinematics, sensor fusion, actuator control, or a robot software stack. The purpose is to show that physical interaction is one grounding category and that it includes both perception and action. \CoALACite{\S 4.2, ``Physical environments''}

\section{Dialogue with humans or other agents}
\textbf{Entry intuition:} Not every environment is a machine or a physical space. Another person can provide instructions, correct a misunderstanding, supply missing information, or respond to a request for help.

\textbf{Coverage:}
\begin{itemize}
  \item Explain dialogue as an external interaction category, not merely a style of generated response.
  \item Cover the paper's examples of accepting instructions and learning from people.
  \item Explain that an agent may ask for help or clarification when it needs information or cannot confidently continue.
  \item Briefly acknowledge interaction among multiple agents, including collaboration, debate, and social simulation, only as categories noted by the paper.
  \item Distinguish a dialogue action from internal reasoning about what someone said: the communication is external grounding; processing the resulting text is internal reasoning.
\end{itemize}

\textbf{Product X example:} The agent may send the on-call engineer a concise question about whether a payment-service deployment or configuration change occurred shortly before the incident. The message is an external dialogue action; the engineer's reply is an observation that can be represented in the agent's working information. The example must not invent a real incident history or imply that the reply establishes causation.

\textbf{Boundary:} Do not turn this section into a chapter on multi-agent architectures, negotiation algorithms, or emotional dialogue systems. \CoALACite{\S 4.2, ``Dialogue with humans or other agents''}

\section{Digital environments: websites, APIs, games, and code execution}
\textbf{Entry intuition:} A request to a website, a game action, or a call to a digital service can affect or inspect a state outside the agent's internal memory, even when no physical robot is involved.

\textbf{Coverage:}
\begin{itemize}
  \item Explain digital environments as external settings that include games, APIs, websites, and general code execution.
  \item Explain the paper's observation that digital grounding is often cheaper and faster than interaction with physical environments or people, making it a convenient testbed.
  \item Explain the paper's treatment of stateless digital APIs---such as search, calculator, or translator APIs---as special ``single-use'' digital environments.
  \item Clarify why a tool-like interface can still be understood as grounding: it connects the agent to external information or computation rather than simply changing the agent's internal working memory.
  \item Avoid implying that every digital environment is stateless. The paper specifically identifies stateless APIs as a special case.
\end{itemize}

\textbf{Product X example:} A call to a payment-service health endpoint, a dashboard query, or a controlled code-execution environment can provide information not already present in working memory. The returned status or result becomes an observation for the agent. A separate external action that changes system state should be presented as another kind of interaction, not conflated with inspection.

\textbf{Planned Mermaid figure 5.2:} \emph{Three external environment categories}. Show physical environments, dialogue with humans or agents, and digital environments as three categories under external grounding. Add concise examples under each category; do not imply the categories are exhaustive of all possible future environments. \CoALACite{\S 4.2, ``Digital environments''}

\section{From perceptual data and digital results to working information}
This section reconnects the three environment categories to the common grounding abstraction.

\textbf{Coverage:}
\begin{itemize}
  \item Compare different input forms: image, audio, tactile signal, natural-language reply, API response, or code-execution result.
  \item Explain that CoALA's account makes textual observations a common interface for the language agent, while acknowledging that other modalities can be supported and that vision-language models can simplify the mapping.
  \item Distinguish the external event from the agent's representation of that event. The representation helps the language model use an observation; it is not the external world itself.
  \item Use one short contrast: a monitoring endpoint may return a structured status object, while a human may reply in natural language; either can be represented as information in working memory.
\end{itemize}

\textbf{Planned Mermaid figure 5.3:} \emph{Different observations, common working interface}. Show three input paths (physical perception, dialogue, digital result) feeding a representation step and then working memory. The diagram is a teaching synthesis of the common textual interface described by CoALA, not a claim that the three paths use an identical conversion algorithm. \CoALACite{\S 4.2; discussion of grounding in \S 4}

\section{A Product X grounding walkthrough}
Use a deliberately limited book-created scenario to demonstrate the concept without moving into the planning chapters.
\begin{enumerate}
  \item The active working information says that the payment endpoint is returning HTTP 500 responses, but does not yet identify the cause.
  \item A digital grounding action queries a health or deployment-status endpoint. The external system returns a result.
  \item The grounding procedure represents the result as an observation in working memory.
  \item If a relevant fact remains unavailable, a dialogue action asks the on-call engineer for clarification. The engineer's response becomes another observation.
  \item The agent's internal reasoning may later interpret these observations, but the reasoning operation is not itself the external grounding action.
\end{enumerate}

\textbf{Planned Mermaid figure 5.4:} \emph{Product X: two grounding interactions}. Show a digital query and a human clarification as separate external interactions, each with a result that may update working memory. Keep the diagram descriptive; do not introduce proposal scoring, selection, or an automatic remediation sequence.

\begin{warningbox}
The Product X sequence is a book-created illustration. It is not a workflow specified by CoALA. Do not present a successful API response, an engineer's reply, or an incident diagnosis as evidence supplied by the paper.
\end{warningbox}

\section{Grounding and the safety boundary of an action space}
The paper notes that external actions can have harmful consequences, just as some internal learning actions can damage the agent's internal state. This section should introduce the concern without presenting a complete safety architecture.

\textbf{Coverage:}
\begin{itemize}
  \item Explain that an external action is not automatically safe simply because it is available to the agent.
  \item Briefly mention the paper's examples of harmful external actions across digital, dialogue, and physical settings.
  \item Note that the paper discusses task-specific safety measures and the need to specify or ablate action spaces in more complex settings.
  \item Distinguish the paper's observation about action-space safety from production engineering policies such as authorization gates, audit logs, rollback controls, and human approvals. Those are implementation topics for Part III, not a complete set of controls prescribed by CoALA.
\end{itemize}
\CoALACite{\S 6, discussion of action-space safety}

\section{Common confusions, recap, and review}
Address these distinctions explicitly:
\begin{itemize}
  \item Grounding is not retrieval: grounding interacts with an external environment; retrieval reads from long-term memory.
  \item An environment is not the same thing as a tool wrapper or a grounding procedure.
  \item An action is not an observation: one is done by the agent; the other is received from the environment.
  \item Dialogue is an external interaction when it communicates with a person or another agent; reasoning about the reply is an internal action.
  \item A digital API can be a special single-use environment, but not every environment is stateless.
  \item Grounding does not decide which action should be proposed, evaluated, selected, or executed; the decision procedure coordinates that work in later chapters.
\end{itemize}

The recap should include a compact comparison of physical, dialogue, and digital environments, and a short flow that links action, environment, feedback, representation, and working memory. Review questions should ask the reader to classify examples and identify the difference between an external action and the observation it produces.

\section*{Source and example note}
The theoretical account is grounded primarily in CoALA Section 4.2. Section 4.6 and Figure 5 are used to position grounding within the complete action space; Section 6 is used only for the paper's brief discussion of action-space safety. Product X is a book-created teaching example. No Fan \& Lan mechanisms or Part III implementation decisions should be imported into this chapter. \CoALACite{\S 4.2, Figure 5, and the action-space safety discussion in \S 6}

\section*{4. Visual and pedagogical plan}
Use three or four Mermaid diagrams, depending on whether the observation-representation figure improves clarity without repeating the grounding loop:
\begin{enumerate}
  \item Grounding interaction: action, environment, feedback, and working memory.
  \item External environment categories: physical, dialogue, and digital.
  \item Different observation forms mapped to a common working-information interface.
  \item Product X digital query and dialogue clarification walkthrough.
\end{enumerate}

All diagram sources will be stored under the following chapter-specific folder:
\begin{center}\small\url{Tex/diagrams/source/part1/chapter05/}\end{center}
Use the shared theme and include \texttt{accTitle} and \texttt{accDescr}. CI will render PNGs before XeLaTeX. The chapter will place figures using the shared \verb|\BookDiagram| macro. No TikZ, ASCII diagrams, or custom LaTeX drawing libraries are permitted.

Use shared teaching environments selectively: a learning objective, a note distinguishing action from observation, a warning about external-action risk, a concise takeaway, chapter recap, and review questions. Keep explanatory prose outside boxes.

\section*{5. Source and attribution plan}
\begin{itemize}
  \item Primary source: Sumers et al., CoALA (TMLR 2024), Section 4.2, ``Grounding actions.''
  \item Architecture context: CoALA Section 4 and Figure 5, to distinguish internal and external action categories.
  \item Physical, dialogue, and digital environment categories: CoALA Section 4.2.
  \item Action-space safety: CoALA Section 6, only for the paper's stated concern and examples.
  \item Use examples only as the CoALA paper presents them or as explicitly marked book-created Product X illustrations.
  \item Trace every theoretical claim to the exact section or figure; do not cite Fan \& Lan in Part I.
\end{itemize}

\section*{6. Explicit exclusions}
Do not re-teach memory or internal actions; Chapters 3 and 4 already cover them. Do not teach proposal, evaluation, selection, or execution control in depth; those topics belong to Chapters 6 and 7. Do not turn the chapter into an API-integration guide, a browser-agent tutorial, a robotics manual, a multimodal model survey, a security architecture, or a comparison of commercial tool frameworks. Do not claim that CoALA prescribes one standard tool protocol, observation schema, or deployment pattern.

\section*{7. Proposed length and review gate}
Estimated depth: approximately 20--28 finished pages, without a hard page limit. The chapter should be long enough to explain all three environment categories, the action/observation boundary, and the Product X example without repeating the overall architecture.

\begin{takeawaybox}[Approval record]
The author approved this outline on 14 November 2026. The chapter draft was prepared from this approved plan, preserving the source map, scope boundaries, examples, and figure requirements. The review checklist below records the matters covered by that approval.

Please review the chapter's purpose, scope, section progression, Product X scenario, Mermaid figure plan, source map, and exclusions. Chapter 5 prose has been drafted from this approved outline.
\end{takeawaybox}
\end{document}
